\chapter{Experiments}
\label{sec:experiments}

\section{Experimental Configuration}
This chapter illustrates how to describe an experimental setup. The values
in Table~\ref{tab:environment} are hypothetical configuration values, not a
record of an experiment performed for this template. A completed thesis should
replace them with the actual settings and explain which choices were fixed
before examining the evaluation results.

\begin{table}[htbp]
  \centering
  \caption{Illustrative configuration; no experiment was executed}
  \label{tab:environment}
  \begin{tabular}{ll}\hline
    Component & Example setting\\\hline
    Candidate set & 100 documents per query\\
    Evaluation cutoff & First 10 results\\
    Configurations & Baseline and reranked variant\\
    Aggregation & Equal weight per query\\
    Output record & One paired score per query\\\hline
  \end{tabular}
\end{table}

\section{Collection and Queries}
The collection description should identify the document source, any filtering
steps and the version used in the evaluation. Query selection deserves the
same level of care. Queries used during development should be distinguished
from those reserved for the final comparison, and any overlap should be
reported rather than left implicit.

Table~\ref{tab:datasets} uses invented group sizes to demonstrate a second table
and its cross-reference. A real report would explain how each group was defined
and why the distinction is relevant to the research question. It would also
state whether group membership was assigned before the system outputs were
examined.

\begin{table}[htbp]
  \centering
  \caption{Invented query groups for the formatting example}
  \label{tab:datasets}
  \begin{tabular}{lrr}\hline
    Group & Development queries & Evaluation queries\\\hline
    Short requests & 20 & 40\\
    Detailed requests & 20 & 40\\
    All requests & 40 & 80\\\hline
  \end{tabular}
\end{table}

% Flush example floats before the continuation-page demonstration.
\clearpage
\section{Execution Procedure}
Run both configurations with the same prepared inputs and retain their complete
ranked outputs. The evaluation stage should read those outputs without altering
the rankings. Separating retrieval from scoring makes it possible to inspect
a result file, correct an evaluation bug and recompute scores without silently
changing the retrieval configuration at the same time.

Each run should have an identifier that connects its configuration, output and
log. A run log can record completion status and unexpected conditions, while
the configuration file records intended settings. If a run fails or produces
an incomplete output, resolve the issue before presenting a comparison as a
complete experiment. The final report should describe consequential deviations
from the original procedure.

\section{Analysis Plan}
The first analysis compares mean query scores. The second examines the
distribution of paired differences, and the third summarises the predefined
query groups. These analyses answer related questions, but none should be used
to imply more than it measures. A group summary based on a small number of
queries may be informative while remaining highly uncertain.

Qualitative inspection can complement these summaries. Select examples using
a stated rule, such as the largest positive and negative changes, and inspect
the retrieved documents for plausible explanations. Such examples illustrate
behaviour; they do not estimate how often that behaviour occurs in the broader
query population. The next chapter uses a small invented result table to
show how these distinctions can be written clearly.
